\subsection{An odd-inertia vanishing criterion}
\label{sec:odd-inertia-criterion}

\begin{proposition}\label{prop:odd-inertia-vanishing}
Let \(R\subseteq L\) be closed normal subgroups of a profinite group
\(G\), and put \(\Gamma=G/R\).
Assume \(C=L/R\) is a central pro-\(2\) subgroup of \(\Gamma\).
Let \(a\in G\) have pro-odd procyclic closure, and assume its closed
normal closure in \(G/L\) is pro-odd.
If \(U\) is a finite continuous \(2\)-primary \(\Gamma\)-module with
\(U^a=0\), then
\[
 H^j_{\rm cont}(\Gamma,U)=0\qquad\text{for every }j\ge0.
\]
\end{proposition}

\begin{proof}
We first show that the closed normal closure \(N\) of the image of
\(a\) in \(\Gamma\) is itself pro-odd.

We use an elementary finite-group fact.
If a finite group \(E_0\) is a central extension of an odd-order group
\(D_0\) by a finite \(2\)-group \(C_0\), then it has a unique odd-order
complement \(J_0\) and \(E_0=C_0\times J_0\).
Indeed \(C_0\) is abelian.  For a central-extension cocycle \(f\),
put
\[
 b(x)=|D_0|^{-1}\sum_{z\in D_0}f(x,z)\in C_0.
\]
Multiplication by \(|D_0|\) is invertible on \(C_0\).
Summing the cocycle identity over \(z\) gives
\(f(x,y)=b(x)+b(y)-b(xy)\), so the extension splits.
Two homomorphic splittings differ by a homomorphism \(D_0\to C_0\),
which is zero by coprimality.
The complement is therefore unique and characteristic, and every
odd-order element of \(E_0\) lies in it.

Now take any finite quotient \(F\) of \(\Gamma\), and let \(C_F\)
be the image of \(C\).  In \(F/C_F\), the image of the assumed
pro-odd normal closure in \(G/L\) is an odd-order normal subgroup.
Its preimage \(E_F\) in \(F\) is a central extension by the \(2\)-group
\(C_F\).  The finite-group fact gives its unique odd-order complement
\(J_F\).  Since \(E_F\) is normal in \(F\), uniqueness makes \(J_F\)
normal in \(F\).
The image of \(a\) has odd order and belongs to \(E_F\), hence to
\(J_F\).  All its \(F\)-conjugates belong to \(J_F\).
Thus the image of \(N\) in every finite quotient of \(\Gamma\) is
odd-order.  Ambient finite quotients are cofinal for the topology
on the closed subgroup \(N\), so \(N\) is pro-odd.

For a pro-odd group and a finite \(2\)-primary continuous module,
averaging in finite quotients gives
\(H^j(N,U)=0\) for every \(j>0\).
Also \(U^N\subseteq U^a=0\), since \(a\) belongs to \(N\).
All rows of the continuous Hochschild--Serre spectral sequence
\[
 H^i(\Gamma/N,H^j(N,U))
       \ \Longrightarrow\ H^{i+j}(\Gamma,U)
\]
are therefore zero.  This proves the assertion, including degree zero.
The spectral sequence and its profinite interpretation are recalled in
\cite[II, Remark~1.35(b), p.~72; Proposition~4.2, p.~87]{MilneCFT}.
\end{proof}

The following application verifies both the central pro-\(2\) layer
and the full normal closure for the exact defining intersection of
\(R_{\rm global}\).
